\chapter{Literature Review}\label{ch:background}

\section{Building Energy Coordination and CityLearn}

CityLearn was designed to standardize reinforcement-learning research for
demand response and urban energy management, and its later version extends the
environment to cover communities that consider both grid needs and occupant
comfort \cite{citylearn2020,citylearnv2}. This is useful here because the
controller must manage several kinds of building information at the same time.
Battery capacity tells us how much energy a building can store and move to a
different time. Heating or cooling demand tells us how much HVAC energy may need
to be controlled. Non-shiftable load is the ordinary electricity use that
cannot be moved by the controller. Recent residential MARL studies also focus
on coordinating many flexible homes while allowing each home to make its own
control decision \cite{charbonnier2025,wilk2024}.

\begin{zhtranslation}
CityLearn 的目标是标准化需求响应和城市能源管理中的强化学习研究，其后续版本进一步
扩展到同时考虑电网需求和住户舒适度的社区控制。这适合本研究，因为控制器需要同时
处理多种建筑信息。电池容量表示一栋建筑能储存多少电，以及能把多少用电转移到其他
时间；供暖或制冷需求表示 HVAC 可能需要控制多少能量；不可转移负荷则是控制器无法
移动的普通用电。近期住宅 MARL 研究也关注如何协调大量具有灵活性的住宅，同时仍让
每栋住宅独立做出控制决定。
\end{zhtranslation}

One control action can affect several results at the same time. Charging or
discharging a battery changes the building's electricity load. Changing HVAC
operation affects both electricity use and indoor temperature. When many homes
act at the same time, their loads add together and can change the community peak
or make the total load rise and fall more quickly. Research on power networks
shows that adding known physical information can help MARL make better control
decisions \cite{physicsMADRL2024}. Research on safe building control also shows
that a small energy error does not automatically mean that comfort and network
limits are satisfied \cite{safeDRL2025,sun2024}. For this reason, the grouping
features in this project describe both how much a building can control and how
much energy it normally needs.

\begin{zhtranslation}
一个控制动作可能同时影响多个结果。电池充电或放电会改变建筑的用电负荷；改变 HVAC
运行会同时影响用电量和室内温度。当很多住宅同时执行动作时，它们的负荷会叠加，从而
改变社区峰值，或者使总负荷上升和下降得更快。电网研究表明，在 MARL 中加入已知的
物理信息能够帮助模型做出更好的控制决定；安全建筑控制研究也表明，能源误差较小并不
代表舒适度和电网限制一定得到满足。因此，本项目的分组特征既描述建筑有多少可控能力，
也描述建筑通常需要多少能量。
\end{zhtranslation}

\section{Grouping Algorithms}

\subsection{K-means}

K-means is used as the original baseline because it is simple and fast
\cite{kmeans1967}. It first places several group centres, assigns every building
to its nearest centre, updates the centres, and repeats these steps. In plain
terms, it tries to put buildings with similar feature values into the same
group. However, its result can change when the starting centres change. It also
works best when groups form simple compact shapes and does not guarantee that
the groups contain similar numbers of buildings. These limitations matter
because an actor trained on a very small group receives less experience than an
actor trained on a large group.

\begin{zhtranslation}
K-means 被用作最初的基线，因为它简单而且计算速度快。它先放置几个组中心，把每栋
建筑分给距离最近的中心，再更新中心并重复这些步骤。简单来说，它尝试把特征数值相近
的建筑放进同一组。但如果起始中心改变，最终结果也可能改变。它更适合形状简单且集中的
分组，而且不能保证每组包含相近数量的建筑。这些限制很重要，因为小组中的 actor 能
获得的训练经验会少于大组中的 actor。
\end{zhtranslation}

\subsection{Gaussian Mixture Model}

A Gaussian mixture model (GMM) treats the data as if they were produced by
several overlapping bell-shaped groups \cite{fraleyraftery2002}. Instead of
only measuring the distance to a centre, GMM estimates the probability that a
building belongs to each group. The building can then be assigned to the group
with the highest probability. This allows GMM to describe groups that have
different widths, directions, or some overlap, so it is more flexible than
K-means. The disadvantage in this experiment is the small dataset: there are
only 25 buildings. A full GMM must estimate many values from these 25 examples,
so its grouping can be unstable or too strongly influenced by a few buildings.

\begin{zhtranslation}
Gaussian mixture model（GMM）把数据看成由几个可能相互重叠、形状类似钟形的组产生
\cite{fraleyraftery2002}。它不只是计算建筑到中心的距离，而是估计每栋建筑属于每一组
的概率，再把建筑分到概率最高的组。这样，GMM 可以表示宽度、方向不同或有一定重叠的
组，因此比 K-means 更灵活。本实验中的缺点是数据量很小，只有 25 栋建筑。完整 GMM
需要从这 25 个样本中估计很多数值，所以结果可能不稳定，也可能受到少数建筑的较大
影响。
\end{zhtranslation}

\subsection{Agglomerative Clustering}

Agglomerative clustering works from small groups to large groups. At the
beginning, every building is its own group. The method then repeatedly merges
the two most similar groups until the required number of groups remains. This
project uses Ward's rule, which chooses the merge that keeps the buildings
inside each new group as similar as possible \cite{ward1963}. The method does
not need random starting centres, so the same input produces the same grouping.
This is useful for a small set of 25 buildings because it removes one source of
randomness from the later MAPPO experiment. Its limitation is that an early
merge cannot be undone, and it does not guarantee equal group sizes.

\begin{zhtranslation}
Agglomerative clustering 按照“从小组到大组”的方式工作。开始时，每栋建筑都是单独
的一组；然后不断合并最相似的两组，直到剩下需要的组数。本项目使用 Ward 规则，选择
能让新组内部建筑尽量相似的合并方式 \cite{ward1963}。该方法不需要随机设置起始中心，
因此相同输入会得到相同分组。对于只有 25 栋建筑的实验，这可以减少后续 MAPPO 实验
中的一个随机来源。它的限制是，前面做出的合并不能撤销，而且不能保证每组大小相等。
\end{zhtranslation}

\subsection{Balanced Spectral Clustering}

Spectral clustering first builds a similarity graph \cite{ng2001spectral}.
Each building is a point in the graph, and two buildings receive a stronger
connection when their selected features are more similar. The method then uses
the main patterns of this graph to create a simpler representation in which
strongly connected buildings are easier to separate from the others. This step
is why the method is called ``spectral'': it studies numerical patterns in the
similarity graph. The buildings are then grouped using this simpler view of
their relationships.

The balanced version used in this project adds one more rule: group sizes must
be equal or differ by at most one building. It creates a fixed number of places
in each group. It then uses a standard matching method, called the Hungarian
assignment algorithm, to place all buildings while keeping the total difference
between buildings and their assigned groups as small as possible. This follows
the idea of balanced clustering
\cite{malinen2014balanced}. This prevents one actor from receiving many
buildings while another actor receives very few. It can also find groups that
are not simple round clusters. However, the result still depends on how
similarity is calculated, and forcing equal sizes may sometimes separate
buildings that would otherwise naturally stay together.

\begin{zhtranslation}
Spectral clustering 首先建立一张相似度图 \cite{ng2001spectral}。图中的每个点代表一栋
建筑；如果两栋建筑所选特征更相似，它们之间的连接就更强。随后，该方法利用图中的
主要结构建立一个更简单的表示，使连接较强的建筑更容易与其他建筑区分开来。它之所以
称为“spectral”，是因为它会分析这张相似度图中的数值规律。最后，再根据这个更简单的
建筑关系进行分组。

本项目使用的 balanced 版本还增加了一条规则：各组大小必须相等，或者最多只相差一栋
建筑。它先为每组设置固定数量的位置，再使用一种称为 Hungarian assignment algorithm
的标准匹配方法分配所有建筑，并尽量减小建筑与其所属组之间的总差异。这与 balanced
clustering 的思想一致 \cite{malinen2014balanced}。
这样可以避免一个 actor 分到很多建筑，而另一个 actor 只分到很少建筑。该方法也能发现
并非简单圆形的组。但它仍然会受到相似度计算方式的影响，而且强制组大小相等有时可能
会把原本适合放在一起的建筑分开。
\end{zhtranslation}

\section{Grouping Features and Redundancy}

Using more features does not always produce better groups. A useful feature
should help distinguish buildings, but it should not simply repeat information
that is already present. This idea is often described as selecting features
that are relevant but not redundant \cite{peng2005mrmr}. For example, mean
heating demand, maximum heating demand, and HVAC rating all describe related
parts of a building's heating requirement. If all three are included, heating
may influence the grouping several times, while battery or ordinary electrical
load is counted only once. The feature ablation therefore tests small
combinations instead of assuming that every available column should be used.

\begin{zhtranslation}
使用更多特征不一定能产生更好的分组。一个有用的特征应当能够区分建筑，但不应该只是
重复已经存在的信息。这通常被称为选择“有用但不重复”的特征 \cite{peng2005mrmr}。
例如，平均供暖需求、最大供暖需求和 HVAC 额定容量都与建筑的供暖需求有关。如果三者
同时加入，供暖因素可能在分组中被重复计算多次，而电池或普通用电只被计算一次。因此，
特征消融实验比较较小的特征组合，而不是假设所有可用特征都应该加入。
\end{zhtranslation}

The selected 3fA set gives one feature to each main part of the task.
\texttt{bes\_capacity\_kwh} tells us how much energy the battery can move.
\texttt{heating\_mean} describes the average heating need during the Vermont
winter. \texttt{nsl\_mean} describes the average ordinary electrical load that
the controller cannot move. These three features are not assumed to be best for
every problem. The experiments in Chapter~\ref{ch:eval} compare different
feature counts and combinations to check whether this simple set works better
than a larger set containing repeated information.

\begin{zhtranslation}
选定的 3fA 为任务中的三个主要部分分别保留一个特征：
\texttt{bes\_capacity\_kwh} 表示电池能够转移多少能量，\texttt{heating\_mean} 表示
Vermont 冬季的平均供暖需求，\texttt{nsl\_mean} 表示控制器无法转移的平均普通用电。
这里并不假设这三个特征对所有问题都最好。第~\ref{ch:eval}章会比较不同的特征数量和
组合，检验这个简单组合是否比包含重复信息的更大特征集合更有效。
\end{zhtranslation}

\section{Soft Routing and Mixtures of Experts}

Mixture-of-experts models use several candidate neural networks and another
small network that chooses between them. The candidate networks are often
called ``experts'', and the choosing network is called a ``router''. In soft
routing, the router gives every candidate a probability instead of immediately
choosing only one. Candidates with higher probabilities have more influence on
the result. Previous Soft MoE and reinforcement-learning studies show that this
kind of selection can help different networks learn different jobs
\cite{puigcerver2024softmoe,obandoceron2024moerl}. In this project, the
candidates are the actors trained for different building groups. For each
building and each time step, the router examines the current building
information and helps select which actor should guide the action.

\begin{zhtranslation}
Mixture-of-experts 模型使用多个候选神经网络，再用另一个小型网络在它们之间进行选择。
候选网络通常称为 experts，负责选择的网络称为 router。在 soft routing 中，router 不会
马上只选择一个候选网络，而是给每个候选网络一个概率；概率越高，它对最终结果的影响
越大。已有 Soft MoE 和强化学习研究表明，这种选择方式可以帮助不同网络学习不同任务
\cite{puigcerver2024softmoe,obandoceron2024moerl}。在本项目中，候选网络就是为不同建筑
组训练的 actors。对每栋建筑和每个时间步，router 会查看建筑当前的信息，并帮助选择
应由哪个 actor 指导动作。
\end{zhtranslation}

The selector is difficult to train when the candidate actors have not learned
useful behaviours. In that case, choosing one actor instead of another provides
little benefit. Training the actors and the selector together from the beginning
can also be unstable because both sides keep changing. The staged experiments
therefore first train actors with fixed groups, then keep those actors unchanged
while training the selector, and finally test whether a small amount of further
actor training is helpful. This means that soft routing is treated as an
improvement built on a good initial grouping, rather than a method that is
expected to create good groups by itself from the beginning.

\begin{zhtranslation}
如果候选 actors 还没有学到有用的行为，这个选择器就很难训练，因为此时选择哪个 actor
都没有明显好处。如果从一开始同时训练 actors 和选择器，也可能不稳定，因为两边都在
不断变化。因此，本报告先使用固定分组训练 actors，再保持这些 actors 不变并训练选择器，
最后测试少量继续训练 actor 是否有帮助。这意味着 soft routing 是建立在良好初始分组
之上的改进方法，而不是期望它从一开始就能自己产生良好分组。
\end{zhtranslation}
